\documentclass[a4paper,12pt]{article}
\usepackage[T2A]{fontenc}
\usepackage[utf8]{inputenc}
\usepackage[russian]{babel}
\usepackage{amsmath,amssymb,mathtools}
\usepackage{geometry}
\geometry{left=20mm,right=20mm,top=20mm,bottom=23mm}
\renewcommand{\familydefault}{\sfdefault}
\usepackage{microtype}
\usepackage{array,booktabs,tabularx}
\usepackage{tikz}
\usepackage{xcolor}
\usepackage{fancyhdr}
\usepackage[unicode=true,colorlinks=true,linkcolor=darkaccent]{hyperref}
\definecolor{accent}{HTML}{087C83}
\definecolor{darkaccent}{HTML}{07595D}
\definecolor{ink}{HTML}{27333B}
\definecolor{wrong}{HTML}{A53639}
\color{ink}
\setlength{\parindent}{0pt}
\setlength{\parskip}{5pt}
\linespread{1.13}
\pagestyle{fancy}
\fancyhf{}
\renewcommand{\headrulewidth}{0.25pt}
\fancyhead[L]{\scriptsize \textcolor{darkaccent}{МАТЕМАТИКА / РЕВЬЮ ДОМАШНЕЙ РАБОТЫ}}
\fancyhead[R]{\scriptsize 10.11.2026}
\fancyfoot[C]{\scriptsize Лёвин Артём Александрович эксклюзивно для Николь Саркисьянц}\fancyfoot[R]{\scriptsize\thepage}
\newcommand{\errorlabel}{\textcolor{wrong}{\textbf{Ошибка.}}}
\newcommand{\keyidea}{\textcolor{darkaccent}{\textbf{Ключевая идея.}}}
\newcommand{\result}{\textcolor{darkaccent}{\textbf{Итог.}}}
\newcommand{\subline}[1]{\par\vspace{3pt}{\color{darkaccent}\bfseries #1}\par}
\begin{document}
\thispagestyle{empty}
\begin{center}
\vspace*{28mm}
{\color{darkaccent}\large МАТЕМАТИКА}\par\vspace{15mm}
{\Huge Ревью домашней работы}\par\vspace{11mm}
{\fontsize{17}{23}\selectfont\color{accent}Математическое ожидание и дисперсия}\par\vspace{15mm}
{\small Подготовка к ЕГЭ по профильной математике}\par
\vfill
{\normalsize Николь Саркисьянц}\par\vspace{8mm}
{\normalsize 10 ноября 2026 г.}\par\vspace{12mm}
{\small Лёвин Артём Александрович}\par
{\small эксклюзивно для Николь Саркисьянц}\par\vspace{12mm}
\begin{tikzpicture}\draw[accent,line width=.9pt] (0,0) .. controls (1.1,.35) and (2.1,-.35) .. (3.2,0) .. controls (4.4,.35) and (5.7,-.35) .. (7.0,0);\end{tikzpicture}
\end{center}
\clearpage
\section*{Сводная таблица}
\textcolor{darkaccent}{Источник:} «Ревью чек-лист 09.11.26.pdf», упражнение к заданию 10.

В таблицу включены только задания, в которых по повторной проверке обнаружены ошибки. Номер задания ведёт к подробному разбору.
\par\par\vspace{4mm}
\small\renewcommand{\arraystretch}{1.35}
\begin{tabularx}{\linewidth}{@{}p{12mm}p{27mm}X p{32mm}p{35mm}@{}}
\toprule
\textbf{№} & \textbf{Тема} & \textbf{Суть ошибки} & \textbf{Ваш ответ} & \textbf{Верный ответ}\\\midrule
\hyperref[task:10]{10} & Дисперсия & Вместо $[M(Y)]^2$ возведена в квадрат вероятность $q$. & \shortstack[l]{$D=11{,}76$\\$\sigma\approx3{,}43$} & \shortstack[l]{$D=5{,}04$\\$\sigma\approx2{,}24$}\\
\bottomrule
\end{tabularx}
\normalsize
\clearpage
\thispagestyle{empty}
\vspace*{\fill}
\begin{center}{\fontsize{25}{30}\selectfont\color{darkaccent}Разбор ошибок}\par\vspace{5mm}
{\large Подробное объяснение и одна задача на закрепление}\end{center}
\vspace*{\fill}
\clearpage
\phantomsection\label{task:10}
\section*{Задание № 10}

\subline{Условие}
Случайная величина $Y$ принимает значения $-1$, $2$ и $5$ с вероятностями $q$, $0{,}4$ и $0{,}6-q$ соответственно. Известно, что $M(Y)=2{,}6$. Найдите $q$, дисперсию и стандартное отклонение.
\subline{Ваше решение}
Из равенства $-q+2\cdot0{,}4+5(0{,}6-q)=2{,}6$ Вы верно получили $q=0{,}2$. Далее правильно рассчитали
\[
M(Y^2)=(-1)^2\cdot0{,}2+2^2\cdot0{,}4+5^2\cdot0{,}4=11{,}8.
\]
Затем записали $[M(Y)]^2=0{,}04$, получили дисперсию $11{,}76$ и стандартное отклонение примерно $3{,}43$.

\textbf{Ваш ответ:} $q=0{,}2$, $D(Y)=11{,}76$, $\sigma(Y)\approx3{,}43$.

\subline{Где именно возникает ошибка}
\textbf{Последний правильный шаг:} $M(Y^2)=11{,}8$.

\errorlabel\ При вычислении $[M(Y)]^2$ Вы использовали $0{,}2^2=0{,}04$ вместо $2{,}6^2=6{,}76$.

\subline{Почему этот шаг неверен}
Число $0{,}2$ является вероятностью $q$. По условию математическое ожидание $M(Y)$ равно $2{,}6$. В формуле дисперсии
\[
D(Y)=M(Y^2)-[M(Y)]^2
\]
возводится в квадрат \emph{математическое ожидание самой случайной величины}, а не вероятность отдельного исхода. Поэтому второе слагаемое равно $2{,}6^2=6{,}76$.

\keyidea\ Различайте три значения: вероятность отдельного исхода $q$, математическое ожидание $M(Y)$ и математическое ожидание квадрата $M(Y^2)$. Каждое выполняет свою роль.

\subline{Исправление с последнего правильного шага}
\[
D(Y)=11{,}8-2{,}6^2=11{,}8-6{,}76=\boxed{5{,}04}.
\]
\[
\sigma(Y)=\sqrt{D(Y)}=\sqrt{5{,}04}=\frac{3\sqrt{14}}5\approx\boxed{2{,}24}.
\]
\subline{Полное правильное решение}
\textbf{1. Найдём неизвестную вероятность.}
\begin{align*}
-q+0{,}4\cdot2+(0{,}6-q)\cdot5&=2{,}6,\\
-q+0{,}8+3-5q&=2{,}6,\\
3{,}8-6q&=2{,}6,\\
-6q&=-1{,}2,\\
q&=0{,}2.
\end{align*}
Третья вероятность равна $0{,}6-0{,}2=0{,}4$. Все вероятности неотрицательны, их сумма равна единице.

\textbf{2. Рассчитаем второй момент.}
\begin{align*}
M(Y^2)&=(-1)^2\cdot0{,}2+2^2\cdot0{,}4+5^2\cdot0{,}4\\
&=0{,}2+1{,}6+10=11{,}8.
\end{align*}
\textbf{3. Найдём дисперсию.}
\[
D(Y)=M(Y^2)-[M(Y)]^2=11{,}8-2{,}6^2=5{,}04.
\]
\textbf{4. Найдём стандартное отклонение.}
\[
\sigma(Y)=\sqrt{5{,}04}=\frac{3\sqrt{14}}5\approx2{,}24.
\]
\subline{Проверка}
\[
0{,}2+0{,}4+0{,}4=1,\qquad
M(Y)=-0{,}2+0{,}8+2=2{,}6.
\]
При расчёте дисперсии через отклонения от среднего:
\[
D(Y)=0{,}2(-1-2{,}6)^2+0{,}4(2-2{,}6)^2+0{,}4(5-2{,}6)^2=5{,}04.
\]
\result\ $q=0{,}2$, дисперсия $5{,}04$, стандартное отклонение $\dfrac{3\sqrt{14}}5\approx2{,}24$.

\subline{Задача на закрепление}
Случайная величина $Z$ принимает значения $-2$, $2$ и $6$ с вероятностями $p$, $0{,}5$ и $0{,}5-p$ соответственно. Известно, что $M(Z)=2$. Найдите $p$, дисперсию и стандартное отклонение $Z$.
\end{document}